\section{Basic Exercises}

For each of the following stochastic matrices $A$, answer the
following items:
\begin{itemize}
\item
  Draw the state diagram for the corresponding Markov chain. You
  should label each node with a positive integer for the corresponding
  column of $A$ (e.g, the node corresponding to the first column
  should be labeled \enquote{1}, to the second column labeled
  \enquote{2}, and so on).
\item
  Determine if $A$ is regular, and write down the smallest $k$ such
  that $A^k$ has strictly positive values. If it is not regular,
  justify your answer.
\item
  Determine the general form solution of the equation
  $(A-I)\vx = \vzero$. You may use a computer to do this,
  but you should express your answer in fractions, not decimals.
\item
  Determine a steady state vector for $A$. If the steady state vector
  is unique, note this. You must do this by hand and show your work.
\end{itemize}
\begin{tasks}[style=enumerate](3)
\task
  \begin{align*}
    \begin{bmatrix}
      0.4 & 0.8 \\ 0.6 & 0.2
    \end{bmatrix}
  \end{align*}
\task
  \begin{align*}
    \begin{bmatrix}
      0.6 & 0 & 0.1 \\
      0.4 & 0.3 & 0 \\
      0 & 0.7 & 0.9
    \end{bmatrix}
  \end{align*}
\task
  \begin{displaymath}
    \begin{bmatrix}
      1 & 1/2 & 1/3 \\
      0 & 1/2 & 1/3 \\
      0 & 0 & 1/3
    \end{bmatrix}
  \end{displaymath}
\task
  \begin{displaymath}
    \begin{bmatrix}
      0.9 & 0.9 \\
      0.1 & 0.1
    \end{bmatrix}
  \end{displaymath}
\task
  \begin{align*}
    \begin{bmatrix}
      0.2 & 0 & 0 \\
      0.4 & 1 & 0.3 \\
      0.4 & 0 & 0.7
    \end{bmatrix}
  \end{align*}
\task
  \begin{align*}
    \begin{bmatrix}
      1 & 0 & 0 \\
      0 & 0 & 1 \\
      0 & 1 & 0
    \end{bmatrix}
  \end{align*}
\task
  \begin{displaymath}
    \begin{bmatrix}
      0.2 & 0 & 0.3 \\
      0.8 & 0.5 & 0 \\
      0 & 0.5 & 0.7
    \end{bmatrix}
  \end{displaymath}
\task
  \begin{align*}
    \begin{bmatrix}
      0 & 0 & 0.3 \\
      1 & 0 & 0.7 \\
      0 & 1 & 0
    \end{bmatrix}
  \end{align*}
\task
  \begin{align*}
    \begin{bmatrix}
      0.9 & 0.6 \\
      0.1 & 0.4
    \end{bmatrix}
  \end{align*}
\task
  \begin{align*}
    \begin{bmatrix}
      1 & 0.3 & 0.4 \\
      0 & 0.7 & 0.1 \\
      0 & 0 & 0.5
    \end{bmatrix}
  \end{align*}
\task
  \begin{align*}
    \begin{bmatrix}
      0 & 0.5 & 0.5 \\
      0.5 & 0 & 0.5 \\
      0.5 & 0.5 & 0
    \end{bmatrix}
  \end{align*}
\task
  \begin{align*}
    \begin{bmatrix}
      0 & 0 & 1 \\
      1 & 0 & 0 \\
      0 & 1 & 0
    \end{bmatrix}
  \end{align*}
\task
  \begin{align*}
    \begin{bmatrix}
      0 & 0 & 0 & 0.5 \\
      1 & 0 & 0 & 0.3 \\
      0 & 1 & 0 & 0.2 \\
      0 & 0 & 1 & 0
    \end{bmatrix}
  \end{align*}
\task
  \begin{align*}
    \begin{bmatrix}
      0.4 & 0.8 & 0.5 & 0 \\
      0.6 & 0.2 & 0 & 0 \\
      0 & 0 & 0 & 0 \\
      0 & 0 & 0.5 & 1
    \end{bmatrix}
  \end{align*}
\end{tasks}

\section{True/False}

\trueFalseHeader

\begin{tasks}[resume, style=enumerate](1)
\task
  For any matrix $A \in \R^{n \times n}$, if $A$ is regular and
  stochastic then $\dim \nul (A - I) = 1$.
\task
  If $A$ is a stochastic matrix, then the all ones vector $\mathbf 1$
  is an eigenvector for $A^T$.
\task
  If a stochastic matrix has a unique steady state vector, then the
  corresponding Markov chain converges to it regardless of starting
  state.
\task
  A stochastic regular matrix $A \in \R^{n \times n}$ has a
  smallest integer $k$ such that $A^k$ has strictly positive
  entries. For such a matrix, $k \leq n$.
\task
  Every stochastic matrix has at least one steady state vector.
\task
  If $A$ and $B$ are stochastic matrices, then so is $AB$.
\task
  Every $A$ stochastic matrix is invertible.
\task
  For any stochastic matrix $A$, if $A$ has a unique stationary state,
  then it is regular.
\task
  If $A$ is a square matrix with strictly positive entries, then there
  is diagonal matrix $D$ such that $AD$ is stochastic.
\end{tasks}

\section{More Difficult Problems}

\begin{tasks}[resume, style=enumerate](1)
\task
  Suppose you've been watching the stock price of your favorite
  company and you've discovered the following trends based on the data
  you've taken:

  \begin{itemize}
  \item
    If the price is \textbf{STEADY} there's a 70\% chance it will
    remain steady the next day, but a 20\% chance it will go
    \textbf{HIGH} and a 10\% chance it will dip \textbf{LOW}.
  \item
    If the price is \textbf{HIGH} there's a 55\% chance it will remain
    high the next day, but a 40\% chance it will go back to being
    \textbf{STEADY}, and just a 5\% chance it will dip all the way
    down to \textbf{LOW}.
  \item
    If the price is \textbf{LOW}, there is a 60\% change it will
    remain low the next day, but a 30\% chance it will reset to
    \textbf{STEADY}, and a 10\% chance it will suddenly jump to
    \textbf{HIGH}.
  \end{itemize}

  Suppose that the price is currently \textbf{STEADY}.  In the long
  term, is more likely to remain \textbf{STEADY}, go up \textbf{HIGH},
  or dip down \textbf{LOW}?  Justify your answer.
\task
  Suppose that we're trying to predict the preformance of the $A_5$
  Soccer Club based on statistics that we've gathered on their recent
  games. We've found that:
  \begin{itemize}
  \item
    After a win, they have a 70\% chance of winning, a 15\% chance of
    drawing, and a 15\% chance losing their next game.
  \item
    After a draw, they have a 20\% chance of winning, a 50\% chance of
    drawing, and a 30\% chance of losing their next game.
  \item
    After a loss, they have a 20\% chance of winning, a 70\% chance of
    drawing, and a 10\% chance of losing their next game.
  \end{itemize}
  \begin{enumerate}
  \item
    If they win their first game, how should we expected them to
    perform overall?  That is, what is the expected percentage of
    wins, draws, and losses in the long term?
  \item
    If they lose their first game, what will their overall record tend
    towards?
  \end{enumerate}
\task
  Consider the following state diagram.
  \begin{center}
    \begin{tikzpicture}
      \node[circle, draw] (1) {1};
      \node[circle, draw] (2) [below left=of 1] {2};
      \node[circle, draw] (3) [below right=of 1] {3};
      \draw[<-] (2) edge["1"] (1);
      \draw[<-] (3) edge["1"] (2);
      \draw[<-] (1) edge["1"] (3);
    \end{tikzpicture}
  \end{center}
  If an edge between two states is not present, the probability of
  transitioning in a single step is $0$.  For example, the probability
  of transitioning from state $3$ to state $2$ in a single step is $0$.
  \begin{enumerate}
  \item
    Determine the transition matrix $T$ for the above diagram. You
    should write $T$ such that the $i$th column of $T$ corresponds to
    the transitions from state $i$. (In the following parts, $T$
    refers to the matrix you determined in this part.)
  \item
    Compute $T^2$, $T^3$ and $T^{2023}$.
  \item
    Determine if $T$ is regular. Justify your answer.
  \item
    Determine if $T$ have a unique steady-state distribution.  Justify
    your answer.
  \end{enumerate}
\task
  Consider the following state diagram.
  \begin{center}
    \begin{tikzpicture}[inner sep=4pt]
      \node[circle, draw] (1) {1};
      \node[circle, draw] (2) [right=of 1] {2};
      \draw[->] (1) edge[bend left, "{$1 - p$}"] (2);
      \draw[->] (1) edge[loop above, "{$p$}"] (1);
      \draw[->] (2) edge[bend left, "{$1 - p$}"] (1);
      \draw[->] (2) edge[loop above, "{$p$}"] (2);
    \end{tikzpicture}
  \end{center}
  \begin{enumerate}
  \item
    Determine the transition matrix $T$ for the above diagram in terms
    of $p$.  In the following parts, $T$ will refer to this matrix.
  \item
    Given $0 < p < 1$, determine $\lim_{k \to \infty} T^k \ve_1$.
    Justify your answer.
  \end{enumerate}
\end{tasks}

\section{Challenge Problems}

\begin{tasks}[resume, style=enumerate](1)
\task
  Demonstrate that any invertible stochastic matrix $A$ with a
  stochastic inverse must be a permutation matrix.  That is, every
  column and every row of $A$ has a single 1, with the rest of the
  entries being 0s.  \textit{Hint:} Consider what can be inferred
  about the entries of a stochastic $A^{-1}$ from its action on $A
  \mathbf{e}_i$ for any elementary basis vector $\mathbf{e}_i$.
\task
 Let $T$ denote the transition matrix of the following state diagram.
  \begin{center}
    \begin{tikzpicture}[inner sep=4pt]
      \node[circle, draw] (1) {1};
      \node[circle, draw] (2) [right=of 1] {2};
      \draw[->] (1) edge[bend left, "{$1 - p$}"] (2);
      \draw[->] (1) edge[loop above, "{$p$}"] (1);
      \draw[->] (2) edge[bend left, "{$1 - p$}"] (1);
      \draw[->] (2) edge[loop above, "{$p$}"] (2);
    \end{tikzpicture}
  \end{center}
  \begin{enumerate}
  \item
    Determine $\lambda_2$, the \textit{second} largest eigenvalue of
    $T$, and a corresponding eigenvector. \textit{Hint:}  $T$ is
    symmetric and, hence, orthogonally diagonalizable.\footnote{The
    second largest eigenvalue tells us about the \textit{rate of
      convergence} to the steady-state distribution. The smaller
    $|\lambda_2|$ is, the faster the convergence to the steady-state
    distribution.}
  \item
    Give a closed-form solution for $T^k [(1-q) \ \ q]^T$ in terms
    of $p$, $q$ and $k$.
  \end{enumerate}
\end{tasks}
